\section{Analytic Design Responses and National Consent}\label{sec:analytic}

The model permits parameter-dependent comparative statics without replacing
its linear-quadratic preferences or the government's sequential agenda.
Write $\alpha=2-\eta$, $\beta=1-\rho$, and $\delta=\rho-\eta$,
with $d=a-c>0$. A firm's $x_i$ is \emph{residual similarity}:
an increase in $x_i$ means \emph{less} differentiation.
All statements in this section concern the canonical
$X_{ij}=(x_i+x_j)/2$ mapping. The strictly interior Nash,
regularity, and positivity conditions are substantive assumptions,
not consequences of writing an interior first-order condition.

\subsection{Differentiation in an international standards union}

For an international union of depth $g$, at a symmetric design
$x_i=y$, set $t=\delta g+\beta y$. Every firm's conditional
quantity is $Q=d/(\alpha+2t)$, and a unilateral design
deviation yields the \emph{Nash} equation
\begin{equation}\label{eq:analytic-is-foc}
 F(y,g;\theta)=\kappa(1-y)-
 \frac{6\beta\alpha d^2}{(\alpha-t)(\alpha+2t)^3}=0.
\end{equation}
This describes an implicitly determined equilibrium design,
not an elementary closed-form $y^*(\theta,g)$.
At fragmentation, the same equation with $g=0$ defines the
symmetric design $z=x^{\SW}$.

Define
\begin{equation}\label{eq:analytic-GK}
 G(t)=\frac{6\beta\alpha d^2(5\alpha-8t)}
 {(\alpha-t)^2(\alpha+2t)^4},\qquad
 K=\kappa-\beta G(t).
\end{equation}

\begin{proposition}[Private-design comparative statics]\label{prop:analytic-A}
Consider a smooth continuation of interior symmetric
\emph{Nash} equilibria satisfying
$\alpha-t>0$, $\alpha+2t>0$, and $K>0$.
Along this continuation,
\begin{equation}\label{eq:analytic-A}
 \frac{\partial y^*}{\partial\kappa}=\frac{1-y}{K}>0,\qquad
 \frac{\partial y^*}{\partial d}
 =-\frac{12\beta\alpha d}
 {(\alpha-t)(\alpha+2t)^3K}<0,\qquad
 \frac{\partial y^*}{\partial g}=\frac{\delta G(t)}{K}.
\end{equation}
If also $5\alpha>8t$, the last derivative has the
sign of $\rho-\eta$: increasing technical compatibility
reduces residual differentiation when $\rho>\eta$ and
increases it when $\rho<\eta$.
\end{proposition}

\begin{proof}
Differentiate the \emph{unilateral} equation
\eqref{eq:analytic-is-foc} implicitly.
It has $F_y=-K$, $F_g=\delta G(t)$ and
$F_\kappa=1-y$.
The derivative with respect to $d$ follows
by differentiating the $d^2$ multiplier.
The sign restrictions are the stated assumptions.
\end{proof}

The firm-scale trade-off depends on the ratio $\kappa/d^2$,
holding the other primitives fixed. Parameter $d$ is the
net demand intercept, not the literal number of consumers.
The proposition is conditional on an interior Nash branch;
an endpoint design may remain fixed as $\kappa$ changes.

\subsection{Different design responses inside and outside a bilateral union}

At the symmetric fragmentation Nash $z$, put
\[
 s_0=\beta z,\quad P=\alpha-s_0,\quad L=\alpha+2s_0.
\]
Let $\ell$ and $h_0$ denote respectively the own and
cross elements of the three-firm design pseudo-gradient
Jacobian, evaluated at fragmentation:
\begin{equation}\label{eq:analytic-lh}
\begin{split}
\ell&=-\kappa+
\frac{6\beta^2d^2(\alpha-2s_0)(2\alpha+s_0)}
 {P^2L^4},\\
h_0&=\frac{3\beta^2d^2(3\alpha-2s_0)}{PL^4}.
\end{split}
\end{equation}
Here $h_0$ is \emph{not} the member--member Cournot
cross coefficient $h$ in Appendix~\ref{app:proofs}.
The direct forcing terms from introducing the $AB$
compatibility edge are $p_0$ (member) and $r_0$ (outsider):
\begin{equation}\label{eq:analytic-pr}
\begin{split}
p_0&=\frac{6\delta\beta d^2(\alpha-2s_0)(2\alpha+s_0)}
 {P^2L^4},\\
r_0&=\frac{6\delta\beta d^2(\alpha^2-2\alpha s_0+4s_0^2)}
 {P^2L^4}.
\end{split}
\end{equation}

\begin{proposition}[Members adjust residual design more than the outsider]
\label{prop:analytic-B}
Suppose the fragmentation design $z\in(0,1)$ is a strict
Nash equilibrium, and the nearby bilateral stationary
solutions are a smooth continuation of genuine interior
Nash equilibria, with uniquely defined member and
outsider designs $u(g),v(g)$. If
\begin{equation}\label{eq:analytic-B-conditions}
 \delta>0,\qquad \alpha>3\beta z,\qquad -\ell>2h_0,
\end{equation}
then
\[
 \boxed{u'(0)>v'(0)>0.}
\]
Thus an infinitesimal increase in bilateral interface
compatibility reduces residual product differentiation
for all firms, but more for member firms.
This is a \emph{local} response at fragmentation,
not a ranking for all positive $g$.
\end{proposition}

\begin{proof}
Taking every firm's \emph{unilateral} derivative before
imposing member symmetry, the implicit-response system is
\begin{equation}\label{eq:analytic-design-system}
 \begin{pmatrix}\ell+h_0&h_0\\2h_0&\ell\end{pmatrix}
 \begin{pmatrix}u'\\v'\end{pmatrix}
 =-\begin{pmatrix}p_0\\r_0\end{pmatrix}.
\end{equation}
The direct-forcing difference is
$p_0-r_0=
6\delta\beta d^2(\alpha-3s_0)/(P^2L^3)>0$.
Under \eqref{eq:analytic-B-conditions}, $h_0>0$ and
$p_0>r_0>0$, while minus the Jacobian in the
full three-firm system is a strictly diagonally
dominant nonsingular M-matrix.
Consequently $u',v'>0$ and, from
\eqref{eq:analytic-design-system},
$u'-v'=(p_0-r_0)/(h_0-\ell)>0$.
The strict Nash-continuation assumption prevents
mistaking a stationary root for a global best response.
\end{proof}

The full-asymmetric Rosen bound
$\kappa>54\beta^2d^2/(\eta^2\alpha^2)$
is a convenient sufficient condition for global
design-Nash uniqueness on the policy cube when
$\delta\geq0$ and $\eta>0$; local interiority still
requires an interior root and continuity.
These conditions, and Proposition~\ref{prop:analytic-B},
are jointly satisfied on the \emph{nonempty open
four-dimensional primitive-parameter region}
independently certified in the supplementary audit.
This does not mean that the result occupies most
of the economically admissible parameter space.

\subsection{Private adjustment and the government's participation margin}

Define the country-$A$ \emph{direct} compatibility effect
at fragmentation, with all firm designs held at $z$:
\begin{equation}\label{eq:analytic-direct}
 D=d^2\left[
 \frac{\rho}{L^2}
 -\frac{2\delta(1+2s_0)}{L^3}
 -\frac{6\delta\alpha}{PL^3}\right].
\end{equation}
Under the conditions of Proposition~\ref{prop:analytic-B},
define the induced national welfare loss
\begin{equation}\label{eq:analytic-loss}
 \begin{split}
 \mathcal L={}&
 \underbrace{\frac{\beta d^2(\eta+2s_0)}{L^3}
          (2u'(0)+v'(0))}_{\text{consumer-surplus loss}}\\
 &+\underbrace{\frac{3\beta d^2(\alpha-2s_0)}{PL^3}
          (u'(0)+v'(0))}_{\text{rivals' design effect on exporter profit}}.
 \end{split}
\end{equation}
\emph{Both terms are added as losses.}
The opposite sign printed for the second term
in the original Stage13 research note was
identified as an error in independent Stage13A
and is not used here.

\begin{proposition}[An analytic local reversal of governmental consent]
\label{prop:analytic-C}
Under the hypotheses of Proposition~\ref{prop:analytic-B},
if the explicitly defined primitive-dependent quantities
satisfy
\begin{equation}\label{eq:analytic-C-conditions}
 0<D<\mathcal L,
\end{equation}
then at $g=0$
\[
 \widehat m'(0)=D>0,\qquad
 m'(0)=D-\mathcal L<0.
\]
Moreover,
$(i-o)'(0)=2m'(0)<0$ and
$(\widehat i-\widehat o)'(0)=2D>0$.
For every qualifying parameter vector,
at \emph{sufficiently small} positive depths
the frozen-design restricted game admits an
international accession accepted by all governments,
whereas the original endogenous-design game rejects
the terminal bilateral union and outsider accession.
The existence of a \emph{globally proposer-optimal}
positive-depth international standard in the frozen game
does \emph{not} follow from this local result alone.
\end{proposition}

\begin{proof}
The home's own-design Nash condition
$\Pi_{A,x_A}=0$ eliminates only that component
of its private-profit response. Direct evaluation
of the domestic consumer-surplus derivative and
the home firm's profit derivatives with respect to
its \emph{rivals'} designs yields
$m'(0)=D-\mathcal L$; see Appendix~\ref{app:analytic-proofs}.
The two outsider-accession slopes follow by permutation
symmetry of the three compatibility edges at fragmentation.
The frozen incumbent accession approval margin also has
positive slope
$(\widehat i-\widehat m)'(0)
=2D+6\delta d^2/(PL^2)>0$,
and $\widehat i'(0)=3D+6\delta d^2/(PL^2)>0$,
so all applicable consent inequalities hold for sufficiently
small positive $g$ by continuity. In the actual design game,
$m(g)-s$ and $i(g)-o(g)$ have strictly negative local
slopes, precluding the positive-depth bilateral path there.
\end{proof}

At the exact witness \eqref{eq:witness},
$D\simeq0.0067679241$ and
$\mathcal L\simeq0.0150697441$, consisting of
a consumer response of approximately $0.005216220$
and a competing-firm profit effect of approximately
$0.009853524$. The independently checked local responses
are $u'(0)\simeq0.060810232$ and
$v'(0)\simeq0.032304817$.
The strict condition \eqref{eq:analytic-C-conditions}
holds not only at this witness but over the
explicit nonzero-volume parameter region
recorded by the exact rational Stage13A certificate.
Unlike the all-depth theorem that follows, this
general-parameter proposition makes only a local claim.
